% GENERATED FILE. Edit canonical lessons, then run npm run generate:books.
% canonical-source-hash: fnv1a64:653783857bdbdccf
% canonical-lessons: SA-C100-kilah, SA-C100-mudgarah, SA-C100-kartari, SA-C100-sucih, SA-C100-sutram

\chapter{Nail, Hammer, Scissors, Needle, Thread}
\label{ch:sa-nail-hammer-scissors-needle-thread}
\hlchaptermodality{100}

\begin{chapteropening}
\textbf{By the end of this chapter:} \emph{I can say a nail, a hammer, scissors, a needle and thread.}

Complete the last lesson of chapter 100: I can say a nail, a hammer, scissors, a needle and thread.
\end{chapteropening}

\section[\texorpdfstring{\sk{कीलः} (kīlaḥ)}{कीलः (kīlaḥ)}]{\sk{कीलः} (kīlaḥ) — a nail (to hammer)}

\label{lesson:SA-C100-kilah}



\begin{quote}
Before the new one: say the Sanskrit for a curtain, then the Sanskrit for a lock.
\end{quote}

\subsection*{You'll want to know: \sk{कीलः}}
\textbf{\sk{कीलः}} — \emph{kīlaḥ} — \textquotedblleft{}a nail (to hammer)\textquotedblright{}.

\begin{grammarlens}[title={What you've built}]
One more word for this chapter.
\end{grammarlens}

\subsection*{Your turn}
\begin{itemize}
\raggedright
  \item \emph{Say it:} \emph{kīlaḥ}
  \item \emph{Say it:} \emph{kīlaḥ}, once more
  \item \emph{Say it:} say \emph{kīlaḥ}
  \item \emph{Recall:} say the Sanskrit for a curtain, then the Sanskrit for a lock, then say \emph{kīlaḥ} again
\end{itemize}

\begin{tcolorbox}[breakable,colback=teal!4,colframe=teal!35!black,title={Before you move on}]
What does \sk{कीलः} mean? (\textquotedblleft{}A nail (to hammer)\textquotedblright{}.) Say it once more.
\end{tcolorbox}
\section[\texorpdfstring{\sk{मुद्गरः} (mudgaraḥ)}{मुद्गरः (mudgaraḥ)}]{\sk{मुद्गरः} (mudgaraḥ) — a hammer}

\label{lesson:SA-C100-mudgarah}



\begin{quote}
Before the new one: say the Sanskrit for a lock, then the Sanskrit for a nail.
\end{quote}

\subsection*{You'll want to know: \sk{मुद्गरः}}
\textbf{\sk{मुद्गरः}} — \emph{mudgaraḥ} — \textquotedblleft{}a hammer\textquotedblright{}.

\begin{grammarlens}[title={What you've built}]
One more word for this chapter.
\end{grammarlens}

\subsection*{Your turn}
\begin{itemize}
\raggedright
  \item \emph{Say it:} \emph{mudgaraḥ}
  \item \emph{Say it:} \emph{mudgaraḥ}, once more
  \item \emph{Say it:} say \emph{mudgaraḥ}
  \item \emph{Recall:} say the Sanskrit for a lock, then the Sanskrit for a nail, then say \emph{mudgaraḥ} again
\end{itemize}

\begin{tcolorbox}[breakable,colback=teal!4,colframe=teal!35!black,title={Before you move on}]
What does \sk{मुद्गरः} mean? (\textquotedblleft{}A hammer\textquotedblright{}.) Say it once more.
\end{tcolorbox}
\section[\texorpdfstring{\sk{कर्तरी} (kartarī)}{कर्तरी (kartarī)}]{\sk{कर्तरी} (kartarī) — scissors}

\label{lesson:SA-C100-kartari}



\begin{quote}
Before the new one: say the Sanskrit for a nail, then the Sanskrit for a hammer.
\end{quote}

\subsection*{You'll want to know: \sk{कर्तरी}}
\textbf{\sk{कर्तरी}} — \emph{kartarī} — \textquotedblleft{}scissors\textquotedblright{}.

\begin{grammarlens}[title={What you've built}]
One more word for this chapter.
\end{grammarlens}

\subsection*{Your turn}
\begin{itemize}
\raggedright
  \item \emph{Say it:} \emph{kartarī}
  \item \emph{Say it:} \emph{kartarī}, once more
  \item \emph{Say it:} say \emph{kartarī}
  \item \emph{Recall:} say the Sanskrit for a nail, then the Sanskrit for a hammer, then say \emph{kartarī} again
\end{itemize}

\begin{tcolorbox}[breakable,colback=teal!4,colframe=teal!35!black,title={Before you move on}]
What does \sk{कर्तरी} mean? (\textquotedblleft{}Scissors\textquotedblright{}.) Say it once more.
\end{tcolorbox}
\section[\texorpdfstring{\sk{सूचिः} (sūciḥ)}{सूचिः (sūciḥ)}]{\sk{सूचिः} (sūciḥ) — a needle}

\label{lesson:SA-C100-sucih}



\begin{quote}
Before the new one: say the Sanskrit for a hammer, then the Sanskrit for scissors.
\end{quote}

\subsection*{You'll want to know: \sk{सूचिः}}
\textbf{\sk{सूचिः}} — \emph{sūciḥ} — \textquotedblleft{}a needle\textquotedblright{}.

\begin{grammarlens}[title={What you've built}]
One more word for this chapter.
\end{grammarlens}

\subsection*{Your turn}
\begin{itemize}
\raggedright
  \item \emph{Say it:} \emph{sūciḥ}
  \item \emph{Say it:} \emph{sūciḥ}, once more
  \item \emph{Say it:} say \emph{sūciḥ}
  \item \emph{Recall:} say the Sanskrit for a hammer, then the Sanskrit for scissors, then say \emph{sūciḥ} again
\end{itemize}

\begin{tcolorbox}[breakable,colback=teal!4,colframe=teal!35!black,title={Before you move on}]
What does \sk{सूचिः} mean? (\textquotedblleft{}A needle\textquotedblright{}.) Say it once more.
\end{tcolorbox}
\section[\texorpdfstring{\sk{सूत्रम्} (sūtram)}{सूत्रम् (sūtram)}]{\sk{सूत्रम्} (sūtram) — thread}

\label{lesson:SA-C100-sutram}



\begin{quote}
Before the new one: say the Sanskrit for scissors, then the Sanskrit for a needle.
\end{quote}

\subsection*{You'll want to know: \sk{सूत्रम्}}
\textbf{\sk{सूत्रम्}} — \emph{sūtram} — \textquotedblleft{}thread\textquotedblright{}.

\begin{grammarlens}[title={What you've built}]
One more word for this chapter.
\end{grammarlens}

\subsection*{Your turn}
\begin{itemize}
\raggedright
  \item \emph{Say it:} \emph{sūtram}
  \item \emph{Say it:} \emph{sūtram}, once more
  \item \emph{Say it:} say \emph{sūtram}
  \item \emph{Recall:} say the Sanskrit for scissors, then the Sanskrit for a needle, then say \emph{sūtram} again
\end{itemize}

\begin{tcolorbox}[breakable,colback=teal!4,colframe=teal!35!black,title={Before you move on}]
What does \sk{सूत्रम्} mean? (\textquotedblleft{}Thread\textquotedblright{}.) Say it once more.
\end{tcolorbox}
